\documentclass[11pt]{article}
\usepackage[a4paper,margin=28mm]{geometry}
\usepackage{amsmath,amssymb,amsthm,mathtools,hyperref}
\newcommand{\ddc}{dd^c}
\newcommand{\dc}{d^c}
\newcommand{\MA}{\operatorname{MA}}
\newcommand{\PSH}{\operatorname{PSH}}
\newtheorem{theorem}{Theorem}
\newtheorem{lemma}[theorem]{Lemma}
\theoremstyle{remark}\newtheorem{remark}[theorem]{Remark}
\title{Exact equations for $L^1$ densities on products with zero lower mass}
\author{}\date{}
\begin{document}\maketitle
\begin{abstract}
Let a compact Hermitian manifold of zero lower Monge--Amp\`ere mass be
multiplied by a normalized projective space. We construct explicit
probability densities in $L^1$ and in no $L^p$, $p>1$, whose individual
thresholds are positive and computed exactly. A genuine unbounded
potential solves the exact finite-mass equation and belongs to the
stretching envelope class; every smaller-background obstacle envelope
has an actual competitor. The threshold calculation retains the
truncations for an arbitrary smooth reference probability. We give a
complete proof of the supplied elliptic-quotient factor construction,
with a direct smooth normalization estimate, and obtain a concrete
fivefold with threshold $240\pi$. The unrestricted zero-lower-mass
$L^1$ equation and potential selection remain open.
\end{abstract}

Use $\ddc=i\partial\bar\partial/(2\pi)$, $\dc=i(\bar\partial-\partial)/(4\pi)$
and ordinary wedge powers. All quasi-psh functions have their genuine
canonical values. The extended measure is the non-pluripolar product
with smooth-background corrections. For a Hermitian form $h$ write
\[
 v_-(h)=\inf_{u\in\PSH(X,h)\cap L^\infty}\int_X(h+\ddc u)^n,
 \qquad b_h(f)=\lim_{j\to\infty}c_h(\min(f,j)).
\]
We use the supplied normalization theorem: $c_h(g)>0$ is the unique
scalar of the normalized continuous bounded-density equation, and
$b_h(f)$ is the maximal scalar of a genuine extended subsolution.
Consequently $b_h$ is decreasing in its density and satisfies
$b_h(tf)=b_h(f)/t$, $t>0$. For bounded densities its truncations are
eventually constant, so $b_h=c_h$. These facts do not identify an
unbounded exact solution.

The stretching class $\mathcal E_{\rm env}(X,h,0)$ consists of genuine
$h$-psh $u\preceq0$ for which every $t\geq1$ admits an $h$-psh $a_t$
with $t^{-1}a_t\leq u$, canonically also at common poles.

\section{The product theorem}
Let $(X_0,h_0)$ be connected and compact of complex dimension $m$,
with $v_-(h_0)=0$, and put $V_0=\int_{X_0}h_0^m>0$.
Let $Y=\mathbb P^k$, $k\geq2$, with normalized Fubini--Study form
$\kappa$, $\int_Y\kappa^k=1$. On $X=X_0\times Y$, of dimension
$N=m+k$, put
\[
 H=h_0+\kappa,\quad C=\binom Nm,\quad
 \mu=\frac{h_0^m}{V_0}\wedge\kappa^k,\quad R=\frac\mu{dV},
\]
where $dV$ is any smooth strictly positive probability on $X$.
All factor forms and functions are pulled back when necessary.
For $0<\alpha<1$ define
\[
 \psi([z])=\log\frac{\sum_{i=1}^k|z_i|^2}{\sum_{i=0}^k|z_i|^2},
 \qquad u_\alpha=1-(1-\psi)^\alpha.
\]
Its unique pole is $o=[1:0:\cdots:0]$. Write
$\MA_\kappa^{\rm ext}(u_\alpha)=q_\alpha\kappa^k$, with density
value $q_\alpha(o)=+\infty$, a harmless measure-zero convention.

\begin{theorem}\label{th:product}
The density $f=R q_\alpha$ is a probability in $L^1(dV)$ and in no
$L^p(dV)$, $p>1$. The genuinely nonclosed product metric satisfies
$v_-(H)=0$, while
\[
 b_H(f)=C V_0>0,\qquad
 \MA_H^{\rm ext}(u_\alpha)=b_H(f)f\,dV.
\]
The displayed solution has supremum zero and finite total mass $CV_0$.
Moreover,
\[
 u_\alpha\in\mathcal E_{\rm env}(X,H,0).
\]
For every $0<\delta<1$
the obstacle envelope $P_{\delta H}(u_\alpha)$ is genuine.
\end{theorem}

\section{The singular density on the closed factor}
The logarithmic norm $\psi$ is a genuine $\kappa$-psh function, because locally
\[
 \kappa+\ddc\psi=\ddc\log\sum_{i\geq1}|z_i|^2\geq0.
\]
The increasing convex function $\chi(t)=1-(1-t)^\alpha$ has
$0<\chi'\leq\alpha$. The composition is $\kappa$-psh, by the smooth
chain rule off its pole and bounded truncation/local regularization
at the pole. Off $o$ it satisfies
$\kappa+\ddc u_\alpha\geq(1-\alpha)\kappa$. It is smooth along
the hyperplane $z_0=0$ and has supremum zero there.

In the affine chart $z_0=1$, put $r=\sum_{i\geq1}|z_i|^2$ and
\[
 F(r)=\log(1+r)+1-\left[1-\log\frac r{1+r}\right]^\alpha,
 \quad Q(r)=rF'(r)
 =\frac{r+\alpha[1-\log(r/(1+r))]^{\alpha-1}}{1+r}.
\]
Radial Stokes with our normalization gives
\[
 \int_{a<|z|^2<b}(\ddc F)^k=Q(b)^k-Q(a)^k.
\]
Indeed the tangential eigenvalue is $F'$ with multiplicity $k-1$ and
the radial eigenvalue is $F'+rF''$; polar integration is the derivative
of $(rF')^k$, with its constant fixed by normalized Fubini--Study mass.
Here $Q(0)=0$ and $Q(\infty)=1$. The omitted hyperplane is smooth and
has zero measure, so the non-pluripolar mass is one. This is a direct
flux calculation, rather than a deduction from zero Lelong number.

Near the pole put $L=1-\log r$. The tangential and radial eigenvalues
are respectively comparable to $r^{-1}L^{\alpha-1}$ and
$r^{-1}L^{\alpha-2}$. Thus
\begin{equation}\label{eq:density}
 q_\alpha(r)\asymp r^{-k}L^{-1-k(1-\alpha)}.
\end{equation}
The radial smooth volume is comparable to $r^{k-1}dr$, proving $L^1$
integrability. For every $p>1$ the substitution $r=e^{1-L}$ leaves
the exponentially increasing factor $e^{k(p-1)L}$, proving divergence.
In particular $q_\alpha\to+\infty$ uniformly as $r\downarrow0$.
Every $q_j=\min(q_\alpha,j)$ is therefore strictly positive and
continuous on all of $Y$, including $o$. Its integral $M_j$ increases
to one. This density convention never alters the canonical pole of $u_\alpha$.

\section{The exact threshold and product measure}
First take smooth $g>0$ on $Y$, with $M_g=\int_Yg\kappa^k$.
Smooth existence in the closed class \cite[Corollary 1]{TW} gives
\[
 A_g=\kappa+\ddc\phi_g>0,\qquad A_g^k=(g/M_g)\kappa^k.
\]
Closed integration forces the source scalar to zero; our potential
is $2\pi$ times the source potential. The pulled-back smooth function
solves the actual full-product equation
\[
 (H+\ddc\phi_g)^N=C h_0^m\wedge A_g^k
                  =(CV_0/M_g)R g\,dV.
\]
Uniqueness of the bounded equation scalar gives $c_H(Rg)=CV_0/M_g$.
For strictly positive continuous $g$, smooth uniform approximation,
its positive lower bound, density monotonicity and reciprocal homogeneity
give the same identity by a multiplicative squeeze. No rough product
limit is needed. Hence $c_H(Rq_j)=CV_0/M_j$.

Put $a=\min_XR>0$ and $A=\max_XR$. The exact original truncation obeys
\[
 R\min(q_\alpha,j/A)\leq\min(Rq_\alpha,j)
                      \leq R\min(q_\alpha,j/a).
\]
Using the continuous-factor identity at both real truncation levels gives
\[
 \frac{CV_0}{M_{j/a}}\leq c_H(\min(f,j))
                       \leq\frac{CV_0}{M_{j/A}}.
\]
Both bounds tend to $CV_0$, proving the exact threshold for arbitrary $dV$.

Now compute the actual unbounded equation separately. Bounded local
products and factor dimensions give
\[
 \MA_H(\max(u_\alpha,-j))
   =C h_0^m\wedge\MA_\kappa(\max(u_\alpha,-j)).
\]
Multiplying by $\mathbf1_{\{u_\alpha>-j\}}$ and taking the defining
non-pluripolar limit yields
\[
 \MA_H^{\rm ext}(u_\alpha)
   =C h_0^m\wedge q_\alpha\kappa^k=CV_0 f\,dV.
\]
This is finite of mass $CV_0$; it is not inferred from scalar convergence.
Smooth bounds on $R$ preserve all the integrability conclusions.

For bounded $h_0$-psh functions $v_j$ with masses tending to zero,
the same local factor expansion gives
\[
 \int_X\MA_H(v_j)=C\int_{X_0}\MA_{h_0}(v_j)\longrightarrow0.
\]
This expansion follows by binomial algebra of local bounded psh products
and smooth-background corrections; no singular current is restricted
to a fiber. Thus $v_-(H)=0$. If $h_0$ were closed positive, bounded
Stokes would fix all its bounded masses to $V_0>0$. Therefore $h_0$
and $H$ are nonclosed.

\section{Actual envelope competitors}
Fix $t\geq1$, choose $\alpha<\beta<1$ and $\varepsilon\beta\leq1$.
The finite constant
\[
 K_t=\sup_{s\geq1}\{ts^\alpha-\varepsilon s^\beta\}-t
\]
makes $a_t=-\varepsilon(1-\psi)^\beta-K_t\leq t u_\alpha$.
Its increasing convex composition has derivative at most
$\varepsilon\beta$, so $\ddc a_t\geq-\kappa$. Its pullback is genuine
$H$-psh, canonically minus infinity at the common pole. Together with
$u_\alpha\leq0$ this proves the precise stretching-class membership.
For $P_{\delta H}$ use $t=1$ and $\varepsilon\beta\leq\delta$.
This supplies an actual $\delta H$-psh competitor below the obstacle.
The upper-regularized envelope remains below the usc obstacle and above
this competitor, hence is genuine and not identically minus infinity.

\appendix
\section{A supplied zero-lower-mass factor, with a complete application proof}
We reproduce the named factor construction supplied in the author's
manuscript; it is not claimed as a new result. Its last normalization
step is proved here directly, without importing a duality formula.

Set $q=e^{2\pi}$, $\Omega=\{[z:w]\in\mathbb P^3:z,w\ne0\}$,
$z,w\in\mathbb C^2$, and $g[z:w]=[z:qw]$.
The free cyclic action has $t=\log(\|w\|^2/\|z\|^2)$ shifted by $4\pi$;
bounded $t$ ranges give proper discontinuity. Its strip $0\leq t\leq4\pi$
is compact: in total-norm-one coordinates both blocks stay away from zero.
It meets every orbit. Thus $X_0=\Omega/\langle g\rangle$ is a connected
compact complex threefold. It maps to $\mathbb P^1\times\mathbb P^1$
with elliptic fiber $\mathbb C/(2\pi\mathbb Z+2\pi i\mathbb Z)$.
Let $\alpha_0,\beta_0$ be the normalized base forms. Then
\[
 \omega=\alpha_0+\beta_0+dt\wedge\dc t>0
\]
descends. The base and vertical kernels have zero intersection.

In the product chart use $x=z_1/z_0$, $y=w_1/w_0$,
$\zeta=\log(w_0/z_0)$, and $\sigma_v=i\,dv\wedge d\bar v/(2\pi)$.
Let $p$ be a smooth probability relative to $\sigma_v$, positive on
$|v|<1$ and zero on $|v|\geq1$. Take a smooth $0\leq\chi_0\leq1$,
supported there and equal to one on $|v|\leq1/2$; put
$\chi=\chi_0(x)\chi_0(y)$ and
\[
 B_\delta=\chi\delta^{-2}p(\zeta/\delta)\sigma_\zeta,
 \quad C_\delta=\sup_{\Theta>0,\ \ddc\Theta=0,\ \int\omega\Theta=1}
                  \int B_\delta\Theta.
\]
Smooth extension by zero makes $B_\delta$ global, and $B_\delta\leq M_\delta\omega$
makes $C_\delta$ finite. The test family is nonempty: direct differentiation
gives $d(\omega^2)=0$, and $\omega^2/(24\pi)$ is a normalized strictly
positive test. In fact $C_\delta\to\infty$. Otherwise Lamari's
lemma \cite[Lemma 3.3]{Lamari}, applied to the smooth possibly nonclosed
$C\omega-B_{\delta_l}$, gives genuine $v_l$ with supremum zero and
$C\omega+\ddc v_l\geq B_{\delta_l}$.

For clarity normalized compactness here follows from local submean and
\cite[I.4.21, I.5.8--5.9]{Dem}. Cover by inner coordinate balls of
radius $r_i$ with outer radius $4r_i$, and choose bounded nonnegative
smooth $\rho_i$ making $v_l+\rho_i$ psh. At a maximum point of $v_l$,
submean on a ball of radius $2r_i$ bounds $\int_{B_i}(-v_l)$.
On every positive-volume overlap this bound supplies a point with
$v_l\geq-A$; submean in the next outer chart bounds its inner ball.
The finite connected overlap graph propagates a global uniform $L^1$
bound. Primary local compactness and finite extraction give a nontrivial
canonical $L^1$ limit $v$. Positivity passes to the limit and
$B_{\delta_l}\rightharpoonup\chi[\zeta=0]$.

On $\chi=1$ the positive closed local current $\ddc(v+\rho)$ dominates
the smooth divisor current and has geometric Lelong number at least one.
Smooth changes of $\rho$ do not change this number. Siu's theorem
\cite[III.8.5]{Dem} gives a proper analytic level set of dimension at most
two containing the divisor segment. This invokes only positive closed
local currents; Demailly's operator is twice ours, with the equivalent
potential-current level equal to two. Its lift contains an open part
of $H_0=\{w_0=z_0\}$. Since $H_0\cap\Omega$ is a projective plane
minus two points, the identity principle gives all of it. Deck invariance
gives $H_k=\{w_0=q^kz_0\}$. Near $[0:1:1:0]$, in the chart $z_1=1$,
these graphs are $z_0=q^{-k}w_0$. Every local defining holomorphic
function vanishes on all of them; for fixed $(w_0,w_1)$ near $(1,0)$,
their distinct values accumulate at zero. The one-variable identity
theorem makes the function vanish on the polydisk, contradicting properness.

Choose $\delta_j\downarrow0$ and smooth positive $\ddc$-closed $\Theta_j$
with $\int\omega\Theta_j=1$ and $J_j=\int B_{\delta_j}\Theta_j>j^4$.
The automorphisms
\[
 T_{a,b,\tau}[z_0:z_1:w_0:w_1]
 =[z_0:z_1+az_0:e^\tau w_0:e^\tau(w_1+bw_0)]
\]
commute with the deck group and translate $(x,y,\zeta)$ by $(a,b,\tau)$.
Average $T^*\Theta_j$ against
$3^{-4}p(a/3)p(b/3)(3\delta_j)^{-2}p(\tau/(3\delta_j))$
relative to $\sigma_a\sigma_b\sigma_\tau$, obtaining $\widehat\Theta_j$.
The parameter set is fixed compact, so $\int\omega\widehat\Theta_j\leq C_1$.
On $V_j=\{|x|,|y|<1/4,|\zeta|<\delta_j/4\}$, every translation to
the support of $B_{\delta_j}$ has $|a|,|b|\leq5/4$,
$|\tau|\leq5\delta_j/4$. The weights are bounded below by a fixed
constant times $\delta_j^{-2}$. The horizontal coefficient
$\widehat F_j$ of $\widehat\Theta_j$ therefore satisfies
$\widehat F_j\geq cJ_j$ on $V_j$.

Put $f_0(a,b)=p(4(a-1))p(b/2)$ and
$w_j(\tau)=4\delta_j^{-2}p(2\tau/\delta_j)$. Average the curves
\[
 [z_0:z_1]\longmapsto
 [z_0:z_1:e^\tau(z_0+uz_1):e^\tau(az_1+bz_0)]
\]
against $w_j(\tau)f_0(a,b)\varepsilon^{-2}p(u/\varepsilon)$,
relative to all four parameter $\sigma$ forms, $0<\varepsilon\leq1/16$.
On these supports $|a-bu|\geq5/8$, so the lines lie in $\Omega$,
and their deck translates are disjoint. Evaluation is a submersion:
variations in $u,a$ span the $w$ directions at $z_1\ne0$, and
variations in $\tau,b$ do so at $z_1=0$. The curve coordinate spans
the remaining direction. Fiber integration of the smooth supported
weights makes the average $T_{j,\varepsilon}$ globally smooth,
semipositive and closed. Compact parameter support bounds its
$\omega$ degree uniformly in $j,\varepsilon$.

At $u=0$, the local equations are $\zeta=\tau$, $y=ax+b$, giving
\[
 T_{j,0}=w_j(\zeta)\sigma_\zeta\wedge K_{f_0},\quad
 K_{f_0}=\int f_0(a,y-ax)\frac{i}{2\pi}
 (dy-a\,dx)\wedge(d\bar y-\bar a\,d\bar x)\,\sigma_a.
\]
On $|x|,|y|\leq1/4$ the weight is bounded below for $a$ in a fixed
disk about one. These covectors have no common nonzero annihilator,
so $K_{f_0}$ has a uniform positive lower bound. For general $u$,
solve $\tau=\zeta-\log(1+ux)$, $b=(1+ux)y-ax$. For each fixed $j$
the resulting coefficients and all their derivatives depend smoothly
on small $u$. Approximate-identity averaging proves local $C^\infty$
convergence to $T_{j,0}$ as $\varepsilon\downarrow0$. No uniform-in-$j$
convergence is asserted.

In the coframe $(x,y,\zeta)$ a positive $(2,2)$ form has cofactor
matrix $Q$, defined by $\xi\wedge\Theta=\operatorname{tr}(H_\xi Q)dV_0$,
$dV_0=\sigma_x\sigma_y\sigma_\zeta$. Its positive root obeys
\[
 Q=2\det(H)H^{-1},\quad H=\sqrt{\det(Q)/2}\,Q^{-1},\quad
 \Theta^{3/2}=\frac3{\sqrt2}\sqrt{\det Q}\,dV_0.
\]
Write $Q_{\widehat\Theta_j}=\left(\begin{smallmatrix}A_j&p_j\\p_j^*&\widehat F_j\end{smallmatrix}\right)$.
On $V_j$, $Q_{T_{j,0}}=\operatorname{diag}(B_j,0)$ with
$B_j\geq c\delta_j^{-2}I_2$. Positivity and the Schur complement give
\[
 \det(Q_{\widehat\Theta_j}+\operatorname{diag}(B_j,0))
 =\widehat F_j\det(B_j+A_j-p_jp_j^*/\widehat F_j)
 \geq\widehat F_j\det B_j.
\]
The root volume on $V_j$ is at least
$c\delta_j^{-2}\sqrt{J_j}\,dV_0$. The smaller region
$|x|,|y|<1/8$, $|\zeta|<\delta_j/8$ has volume comparable to
$\delta_j^2$. For this fixed $j$, choose $\varepsilon_j>0$ so that
local smooth convergence preserves half this root integral. Then
$\Gamma_j=\widehat\Theta_j+T_{j,\varepsilon_j}$ is smooth, positive,
$\ddc$-closed, has bounded $\omega$ degree and root volume
$R_j\geq c\sqrt{J_j}/2\to\infty$.

For any Hermitian $h_0\leq A\omega$, take $\eta_j^2=\Gamma_j$,
$P_j=\int h_0\eta_j^2\leq C$, and $\mu_j=\eta_j^3/R_j$.
General Hermitian smooth existence \cite[Corollary 1]{TW} gives
$h_j=h_0+\ddc v_j>0$, $h_j^3=c_j\mu_j$, $\sup v_j=0$.
The source scalar is not set to zero here. Smooth Stokes and matrix
AM--GM give
\[
 P_j=\int h_j\eta_j^2\geq c_j^{1/3}R_j^{2/3},\qquad
 0<\int h_j^3=c_j\leq P_j^3/R_j^2\longrightarrow0.
\]
These actual smooth potentials prove $v_-(h_0)=0$ for every $h_0$.

On a fiber $\zeta=s+i\tau$ has both periods $2\pi$ and $t=2s$ up to
a constant. Thus $dt\wedge\dc t=ds\wedge d\tau/\pi$ has fiber
area $4\pi$. Its square is zero and $(\alpha_0+\beta_0)^3=0$, whence
\[
 \int_{X_0}\omega^3
 =3\int(\alpha_0+\beta_0)^2\,\cdot4\pi=24\pi.
\]
Furthermore $\ddc t=\beta_0-\alpha_0$, so
$d\omega=-dt\wedge(\beta_0-\alpha_0)\ne0$.
Taking $h_0=\omega$, $Y=\mathbb P^2$ in Theorem~\ref{th:product}
gives the concrete fivefold and $b_H(f)=\binom53\,24\pi=240\pi$.

\begin{remark}
The sufficient geometric condition is the split background and displayed
density family. No positive lower mass, $L^p$ hypothesis, dual maximizer
or moment bound is assumed. We prove exactness by factor calculation,
rather than selecting a limit of exponential solutions. This does not
resolve arbitrary $L^1$ data at zero lower mass, uniqueness of the
unbounded solution, a greatest-member identity or literature priority.
The named factor is supplied prior work; only the product equation and
its exact constants are the candidate contribution.
\end{remark}
\begin{thebibliography}{9}
\bibitem{TW} V. Tosatti and B. Weinkove, \emph{The complex Monge--Amp\`ere
equation on compact Hermitian manifolds}, arXiv:0910.1390v2 (2009),
Corollary 1 and closed normalization, pp.~1--2;
\url{https://arxiv.org/pdf/0910.1390v2}.
\bibitem{Lamari} A. Lamari, \emph{Courants k\"ahl\'eriens et surfaces
compactes}, Ann. Inst. Fourier \textbf{49} (1999), 263--285,
Lemma 3.3 and proof, pp.~275--276;
\url{https://www.numdam.org/article/AIF_1999__49_1_263_0.pdf}.
\bibitem{Dem} J.-P. Demailly, \emph{Complex Analytic and Differential
Geometry}, 21 June 2012, I.4.21 and I.5.8--5.9, pp.~36--40;
III.5.7--5.9, p.~159; III.8.5 and complete relevant proof, pp.~174--180;
\url{https://people.math.harvard.edu/~demarco/Math274/Demailly_ComplexAnalyticDiffGeom.pdf}.
\end{thebibliography}
\end{document}
